\section{Separating Conditional Noise from Viewing-Law Sensitivity}
\label{app:viewvariance}
This post-outcome follow-up studies all four nonzero ranged-amplitude contrasts
on 10028 and 10049. The zero 10076 contrasts and fixed-amplitude outcomes remain
in Appendix~\ref{app:momentdevelopment}; this is not a new three-stack positive
result. The following construction combines classical conditional covariance,
concentration and distributional robustness
\citep{janon2014sobol,maurer2009bernstein,duchi2019variance}. Its numerical
comparison does not establish novelty or experimental calibration.

\subsection{A conditional-noise envelope and its assumptions}
Partition the amplitude interval into cells $I_j$. For each view $R$, simulate
$L$ independent noises and set
\begin{align*}
 Z_{\ell j}&=\mathbf1\{\max_{a\in I_j}f(a m(R)+\epsilon_\ell)>c\},\\
 X&=\max_j L^{-1}\textstyle\sum_\ell Z_{\ell j},\qquad
 T=L^{-1}\textstyle\sum_\ell\max_j Z_{\ell j}.
\end{align*}
Here $X\le T$ pointwise. For each fixed amplitude, its cell's empirical
frequency is bounded by $X$. Consequently
$\eta_L(R):=\mathbb E[X\mid R]$ dominates every fixed-amplitude event
probability. It also dominates a random physical amplitude's probability if
\emph{amplitude and measurement noise are conditionally independent given
view}. This qualifier is essential. The older $T$ construction protects the
larger class of amplitudes that may depend on realized noise, when the specified
noise marginal and particle independence are retained. Grouping thus changes
this part of the protected null. Image-fitted scales or image-dependent particle
selection do not automatically satisfy the stronger premise. The finite-$L$
maximum bias is retained in $\eta_L$; it is not the exact amplitude-profiled
event probability.

Take two independent noise groups at each view and a center $b\in[0,1]$ fixed
by an older independent simulation. Conditional independence gives
\[
 \mathbb E[(X_1-b)(X_2-b)\mid R]=(\eta_L(R)-b)^2.
\]
Write $\mu=\mathbb E_Q\eta_L$ and $V=\operatorname{Var}_Q\eta_L$.
For $w=dP/dQ\le\kappa$ with $\mathbb E_Qw=1$, Cauchy--Schwarz gives
\[
 \mathbb E_P\eta_L\le\mu+\sqrt{(\kappa-1)V},
 \qquad \mathbb E_P\eta_L\le\kappa\mu.
\]
The first inequality alone also holds with $\kappa-1$ replaced by a supplied
$\chi^2(P\Vert Q)$ upper bound. Under only that weaker assumption, the
$\kappa\mu$ branch must be discarded. Neither divergence radius is estimated
from experimental particles here.

\paragraph{Finite-simulation bound.}
Apply the established empirical Bernstein inequality to independent \emph{views},
not their correlated replicas. With $M$ observations in $[l,u]$ and unbiased
sample variance $s^2$, its one-sided radius at failure probability $d$ is
\[
 r_d=\sqrt{2s^2\log(2/d)/M}
       +\frac{7(u-l)\log(2/d)}{3(M-1)}.
\]
Allocate $\delta/4$ to each tail of the mean of
$\overline X=(X_1+X_2)/2$, yielding $[L_\mu,U_\mu]$, and $\delta/2$ to an
upper mean $U_Y$ of $Y=(X_1-b)(X_2-b)$. Its exact range is
$[-b(1-b),\max(b^2,(1-b)^2)]$. Define
\begin{align*}
 V_U&=\min\{1/4,\max[0,U_Y-\operatorname{dist}(b,[L_\mu,U_\mu])^2]\},\\
 p_0&=\min\{1,\kappa U_\mu,U_\mu+\sqrt{(\kappa-1)V_U}\}.
\end{align*}
Except on one event of probability at most $\delta$, both $\mu\le U_\mu$
and $V\le V_U$, since $\mathbb E Y=V+(\mu-b)^2$. The minimum is taken
on that same event. Conditional on calibration, independent test indicators
are then dominated by independent Bernoulli variables of parameter $p_0$.
The binomial rule of Appendix~\ref{app:momentdevelopment} at level
$\alpha-\delta$ therefore has joint false-rejection probability at most
$\alpha$, in exact arithmetic and under the stated simulator. This is a
conditional-model composition of known results, not a new concentration theorem.

\subsection{Matched draws, classical comparators and concentration slack}
For each stack, the follow-up uses 32,768 new Haar calibration views with two
groups of $L=32$ noises, 16 amplitude cells, and 131,072 independent Haar
held-out views. Both maps and both contrasts share draws within each stage.
The three initially declared methods use the individual envelope, grouped
envelope and grouped viewing variance. They have the same projection and noise
budgets; this is not an optimized allocation comparison with the earlier
131,072-view Clopper--Pearson procedure. All 360 projections are retained.

A focused independent model audit requested classical CVaR comparators. Under
the density cap, the exact population worst case is the upper tail mean
\citep{rockafellar2002cvar}:
\[
 \sup_{0\le w\le\kappa,\ \mathbb E_Qw=1}\mathbb E_Q[w\eta_L]
 =\min_{0\le t\le1}\{t+\kappa\mathbb E_Q(\eta_L-t)_+\}.
\]
An extremal weight fills the upper tail, using fractional weight at a boundary
atom when needed. Conditional Jensen permits replacing $\eta_L$ by
$\overline X$ to obtain an upper bound. We implement two separate confidence
procedures: a uniform DKW bound on the empirical stop-loss curve, and an
empirical-Bernstein bound at a threshold chosen on half the views and evaluated
on the other half. For DKW, add
$\sqrt{\log(2/\delta)/(2M)}(1-t)$ to the empirical stop loss before minimizing.
We also implement a mean-only variance envelope and an unpaired
$\operatorname{Var}(\overline X)$ envelope. Each method has its own
$\delta=.001$ budget; no unadjusted cross-method minimum is used. All 480
additional projections and paired critical-count differences are retained.

At $\kappa=1.1$, the removed-candidate probability bounds for 10049 are
.43095/.44761 (power/combined) using paired variance, versus .43614/.45220
using split CVaR, .44104/.45722 using DKW CVaR, and .44472/.46102 using
unpaired variance. These are comparisons at one simulation budget, not a
contradiction of population CVaR optimality. The CVaR relaxation retains
conditional-noise variability in $\overline X$.

The variance upper bounds are predominantly concentration slack. For the
removed 10049 power/combined cases, the signed centered-product estimates are
$-9.84\times10^{-5}$/$-3.23\times10^{-5}$, the square-root terms are
$1.71\times10^{-4}$/$1.73\times10^{-4}$, and the range terms are
$3.40\times10^{-4}$/$3.32\times10^{-4}$. Distance subtraction is zero;
final $V_U$ values are $4.13\times10^{-4}$/$4.72\times10^{-4}$. This does
not measure zero population viewing variance. No independent calibration-repeat
or optimal-replication comparison has been completed.

At $n=10^4$, $\kappa=1.1$, the paired method's Haar-alternative rejection
projections are .116 [.041,.256] and .185 [.075,.362] on 10049, versus
.00108 and .00816 on 10028. At $n=10^5$, 10049 projections rise to
.443 [.032,.942] and .773 [.167,.993]. These wide intervals transform
pointwise single-image probability intervals; they do not assess calibration
variability or a paired between-method effect. All tests remain weak at
$\kappa=2$ on the amplitude-one Haar alternative.

\subsection{Fresh preferential views and amplitude controls}
We additionally generate all nine laws obtained by conditioning Haar rotations
on $|R_{3j}|\ge1-1/\kappa$, for the three object axes and
$\kappa\in\{1.1,2,5\}$. Each law has density exactly $\kappa$ on its
retained cap union. Complete rotations retain random in-plane angles.
Each stack/law has 65,536 new independent views; all true/removed maps,
amplitudes .9/1/1.1 and both contrasts are evaluated with common noise.
There are 1,179,648 independent rotations across the 18 stack/law runs.
The viewing law is supplied, not inferred from images. These caps are not the
extremal event-probability tilt.

For combined moments on 10049 at $n=10^4$, $\kappa=1.1$, true-map rejection
against the removed candidate ranges across axes from .00357--.00579 at
amplitude .9, .312--.386 at amplitude one, and .984--.988 at amplitude 1.1.
At $\kappa=2$, the amplitude-one range is .0000605--.00774; at $\kappa=5$
it is below $2.1\times10^{-11}$. These ranges summarize point projections,
not confidence intervals over viewing laws. All 1,944 projections with
pointwise intervals are archived. Amplitude sensitivity prevents using the
favorable amplitude-1.1 result as a general detection claim.

Partitioning each law's independent views into 64 disjoint groups of 1,024
particles gives actual repeated-group outcomes. Across 648 correct-null
method/law/amplitude/score cells, the largest rejected count is 2/64. The
pointwise exact interval is wide, and cells share draws; these cannot be pooled
as independent repetitions. Calibration remains fixed, so this is not an
empirical unconditional error study. At $n=10^4$, the largest correct-null
point projection is .000259. Thirty-six first-view physical sums and 216
direct scores agree within $1.1\times10^{-11}$, and every count/critical value
is replayed. Frequencies are nonzero, unique and lack antipodal pairs; the
first transfer profile matches its archived source exactly.

An additional verifier selects 328 of the 32,768 original calibration views
per stack with a separate seed. Separable physical-cell sums, independently
written moment polynomials, and companion-matrix stationary roots reproduce
all selected grouped and individual counts: 167,936 cubics and 10,496 group
comparisons. This is a numerical check on about 1\% of the calibration views,
not a global rounding certificate. Oracle scores, full-region removal, one
known transfer profile and specified Gaussian noise remain substantial limits.
The focused audit supports the conditional algebra but supplies no acceptance
assessment; the three full-paper rejections remain unresolved.
